\addcontentsline{toc}{subsubsection}{Solution 18. The Geometric-Mean Trap}

\begin{dlsummary}[title={\bfseries Solution 18. The Geometric-Mean Trap},breakable]

\begin{proof}
Write $P(x,y)$ for the inequality in the statement. It is symmetric
in $x$ and $y$. On the diagonal, $P(x,x)$ is equivalent to
\[
\frac{2x}{f(x)}\leq\frac{2f(x)}x,
\]
which, since $x,f(x)>0$, is equivalent to
\[
f(x)\geq x. \tag{1}
\]

\medskip
\noindent
\textbf{Step 1. Every admissible pair is diagonal.}

Suppose $P(x,y)$ holds for some $x\ne y$. If $f(x)\geq x$, then
$P(x,x)$ also holds, contradicting uniqueness for this fixed $x$.
Thus $f(x)<x$. By symmetry, $f(y)<y$ as well. Therefore,
\[
\frac{x}{f(y)}+\frac{y}{f(x)}
>\frac xy+\frac yx\geq2,
\]
whereas
\[
\frac{f(x)}x+\frac{f(y)}y<2.
\]
This contradicts $P(x,y)$. Hence every admissible pair has $x=y$.

For each $x>0$, an admissible $y$ exists, so it must equal $x$.
Consequently, $P(x,x)$ holds for every $x>0$, and (1) gives
\[
f(x)\geq x\qquad(x>0). \tag{2}
\]

\medskip
\noindent
\textbf{Step 2. Use a geometric-mean choice.}

Suppose $f(a)>a$ for some $a>0$, and put
\[
u=\frac{f(a)}a>1,
\qquad y=a\sqrt u=\sqrt{a f(a)}.
\]
By (2), $f(y)\geq y$. Hence
\[
\frac{a}{f(y)}+\frac{y}{f(a)}
\leq\frac ay+\frac{y}{f(a)}
=\frac1{\sqrt u}+\frac1{\sqrt u}
=\frac2{\sqrt u}<2.
\]
On the other hand,
\[
\frac{f(a)}a+\frac{f(y)}y\geq u+1>2.
\]
Thus $P(a,y)$ holds. Since $u>1$, $y\ne a$, contradicting Step 1.
Therefore $f(a)\leq a$ for every $a>0$. Together with (2), this
proves $f(x)=x$ for all $x>0$.

\medskip
\noindent
\textbf{Verification.}

For $f(x)=x$, the inequality becomes
\[
\frac xy+\frac yx\leq2.
\]
By AM--GM, the left-hand side is at least $2$, with equality if and
only if $x=y$. Thus exactly one $y>0$ works for each $x>0$, namely
$y=x$.
\end{proof}

\end{dlsummary}